\documentclass[11pt]{article}
\usepackage[margin=1in]{geometry}
\usepackage{amsmath, amssymb, amsthm}
\usepackage{booktabs}

\title{Problem 464: M\"obius Function and Intervals}
\author{}
\date{}

\newtheorem{theorem}{Theorem}
\newtheorem{lemma}[theorem]{Lemma}

\begin{document}
\maketitle

\section{Problem Statement}

Let $P(a,b)$ and $N(a,b)$ count values of $n \in [a,b]$ with $\mu(n) = 1$ and $\mu(n) = -1$, respectively. Define $C(n)$ as the number of pairs $(a,b)$ with $1 \le a \le b \le n$ satisfying
\[
99\,N(a,b) \le 100\,P(a,b) \quad\text{and}\quad 99\,P(a,b) \le 100\,N(a,b).
\]
Given: $C(10) = 13$, $C(500) = 16676$, $C(10000) = 20155319$. Find $C(20\,000\,000)$.

\section{Mathematical Foundation}

\begin{theorem}[Reduction to 2D dominance]
\label{thm:dom}
Define $\mathrm{prefP}(i) = |\{n \le i : \mu(n) = 1\}|$, $\mathrm{prefN}(i) = |\{n \le i : \mu(n) = -1\}|$, and
\[
X(i) = 100\,\mathrm{prefP}(i) - 99\,\mathrm{prefN}(i), \qquad Y(i) = 100\,\mathrm{prefN}(i) - 99\,\mathrm{prefP}(i).
\]
Then $C(n) = |\{(j, b) : 0 \le j < b \le n,\; X(b) \ge X(j),\; Y(b) \ge Y(j)\}|$.
\end{theorem}

\begin{proof}
Substitute $P(a,b) = \mathrm{prefP}(b) - \mathrm{prefP}(a-1)$ and the analogous formula for~$N$ into both inequalities. The first inequality rearranges to $X(b) \ge X(a-1)$ and the second to $Y(b) \ge Y(a-1)$. Setting $j = a - 1$ gives the stated 2D dominance formulation.
\end{proof}

\begin{lemma}[Sum identity]
\label{lem:sum}
$X(i) + Y(i) = \mathrm{prefP}(i) + \mathrm{prefN}(i)$, the count of squarefree integers up to~$i$.
\end{lemma}

\begin{proof}
Direct computation: $X(i) + Y(i) = (100 - 99)(\mathrm{prefP}(i) + \mathrm{prefN}(i)) = \mathrm{prefP}(i) + \mathrm{prefN}(i)$. This is the number of squarefree integers $\le i$, which is non-decreasing.
\end{proof}

\begin{theorem}[CDQ divide-and-conquer]
\label{thm:cdq}
The 2D dominance count with ordering constraint $j < b$ can be computed in $O(n \log^2 n)$ time.
\end{theorem}

\begin{proof}
Divide the index range $[0, n]$ at the midpoint. Recursively count same-half pairs. For cross-pairs ($j$ in the left half, $b$ in the right):
\begin{enumerate}
    \item Sort both halves by $X$-value.
    \item Two-pointer sweep: for each $b$ in increasing $X$ order, insert all $j$ from the left half with $X(j) \le X(b)$ into a Fenwick tree keyed by coordinate-compressed $Y$-values.
    \item Query the Fenwick tree for the prefix sum up to $Y(b)$.
\end{enumerate}
Each level processes $O(n)$ elements with $O(\log n)$ overhead from sorting and Fenwick operations. With $O(\log n)$ levels, the total is $O(n \log^2 n)$. Every valid pair $(j, b)$ with $j < b$ is counted exactly once: either in the same recursive half, or as a cross-pair in the merge step.
\end{proof}

\section{Algorithm}

\begin{enumerate}
    \item Compute $\mu(1), \ldots, \mu(n)$ via a linear sieve in $O(n)$.
    \item Build prefix sums $\mathrm{prefP}$, $\mathrm{prefN}$ and coordinates $X$, $Y$.
    \item Apply CDQ divide-and-conquer to count valid pairs.
\end{enumerate}

\section{Complexity Analysis}

\begin{itemize}
    \item \textbf{Time:} $O(n \log^2 n)$. The M\"obius sieve is $O(n)$; the CDQ step dominates.
    \item \textbf{Space:} $O(n)$ for the sieve, prefix sums, and Fenwick tree.
\end{itemize}

\section{Verification}

\begin{center}
\begin{tabular}{cr}
\toprule
$n$ & $C(n)$ \\
\midrule
10 & 13 \\
500 & 16\,676 \\
10\,000 & 20\,155\,319 \\
\bottomrule
\end{tabular}
\end{center}

\section{Answer}

\[
\boxed{198775297232878}
\]

\end{document}
